\section*{Capítulo \ref{ch:b1:hotspots} --- Pontos quentes e vulcanismo intraplaca}

\begin{solution}{exo:b1:hotspots:1}
(a) $4205 + 5274 = \qty{9479}{m}$, quase \qty{9.5}{km}. O vulcão é mais pesado que a água que substitui: a placa dobra para baixo sob sua carga, de modo que sua base fica mais profunda que o fundo distante, e sua altura real é maior. (b) $74000/0.1 = \num{740000}$ anos com o suprimento do Kilauea; $74000/0.025 = \num{3.0e6}$ anos à taxa média da cadeia. (c)
$0.12\times19 = \qty{2.3}{km^3}$, cerca de $\num{3e-5}$ do Mauna Loa: três milésimos de um por cento.
\end{solution}

\begin{solution}{exo:b1:hotspots:2}
(a) $0.92\times\pi/180 = \qty{0.0161}{rad.Myr^{-1}}$. (b) Em radianos por milhão de anos vezes quilômetros, $\omega R\sin\Delta$ é em quilômetros por milhão de anos, isto é, milímetros por ano: $0.016057\times6371 = 102.3$, logo $v = 102.3\sin\Delta$, o que dá \qty{51.1}{mm.yr^{-1}} a \ang{30}, \qty{88.6}{mm.yr^{-1}} a \ang{60}, \qty{102.3}{mm.yr^{-1}} a \ang{90} e \qty{97.3}{mm.yr^{-1}} a \ang{108.0}. (c) A velocidade é proporcional à distância $R\sin\Delta$ ao eixo, maior sobre o grande círculo a \ang{90} do polo; além dele o ponto se aproxima do eixo de novo, do lado do polo oposto.
\end{solution}

\begin{solution}{exo:b1:hotspots:3}
(a) Decão: $\num{e6}/\num{5e5} = \qty{2}{km}$. Atlântico Central: $\num{2e6}/\num{7e6} = \qty{0.29}{km}$, cerca de \qty{290}{m}. (b) A média do Decão fica perto de sua seção de \qty{2.5}{km}: a província é um platô espesso e contínuo. (c) A área do Atlântico Central é a de toda a região, em quatro continentes, sobre a qual seus restos são encontrados: lavas, diques e soleiras hoje separados por um oceano e muito erodidos. A rocha se concentra em alguns lugares e está ausente da maior parte da área, de modo que a média não diz nada sobre uma única pilha.
\end{solution}

\begin{solution}{exo:b1:hotspots:4}
(a) $\rho_m\alpha\,\Delta T = 3300\times\num{3.2e-5}\times250 =
\qty{26.4}{kg.m^{-3}}$. (b) A coluna da elevação tem a massa extra $(\rho_m - \rho_w)h = 2300\times690 = \qty{1.587e6}{kg.m^{-2}}$, que a camada leve deve equilibrar: $H = \num{1.587e6}/26.4 = \qty{6.0e4}{m}$, cerca de \qty{60}{km} de rocha quente.
\end{solution}

\begin{solution}{exo:b1:hotspots:5}
(a) Em A ($\lambda = 0$, $\varphi = 0$), $\vec e = (0, 1, 0)$ e
$\vec n = (0, 0, 1)$; com $\sin\ang{270} = -1$ e $\cos\ang{270} = 0$, $\vec t_A = -\vec e = (0, -1, 0)$. Em B ($\lambda = \ang{90}$), $\vec e = (-1, 0, 0)$, $\vec n = (0, 0, 1)$ e $\vec t_B = (1, 0, 0)$. (b) $\vec t_A\times\vec t_B = (0, 0, 1)$. (c) Com $\vec p' = (0, 0, 1)$ e $\vec r_A = (1, 0, 0)$, $\vec p'\times\vec r_A = (0, 1, 0)$, e $(\vec p'\times\vec r_A)\cdot\vec t_A = -1 < 0$: o polo é $-\vec p' = (0, 0, -1)$, o polo Sul. Um movimento para oeste ao longo do equador é horário visto de acima do polo Norte, logo anti-horário visto de acima do polo Sul. (d) A \ang{90} do polo, $\omega = 50/6371 =
\qty{7.85e-3}{rad.Myr^{-1}} = \qty{0.45}{\degree\per\Myr}$.
\end{solution}

\begin{solution}{exo:b1:hotspots:6}
(a) $\sum x_i v_i = 90000 + 240000 + 342000 = \num{672000}$ e $\sum x_i^2 = \num{7.0e7}$, logo $\omega = \qty{0.0096}{rad.Myr^{-1}}$, ou \qty{0.55}{\degree\per\Myr}. (b) As previsões são 28.8, 48.0 e \qty{57.6}{mm.yr^{-1}}; os resíduos $+1.2$, $0.0$ e $-0.6$
\unit{mm.yr^{-1}}. (c) Numa rotação rígida a velocidade é proporcional à distância ao eixo, e um ponto sobre o eixo não se move: não há intercepto a ajustar.
\end{solution}

\begin{solution}{exo:b1:hotspots:7}
(a) $v = \qty{50e-3}{m}/\qty{3.156e7}{s} = \qty{1.58e-9}{m.s^{-1}}$, logo
$B = 2300\times\num{3e8}\times\num{1.58e-9} = \qty{1090}{kg.s^{-1}}$ e $Q = (910/\num{3.2e-5})\times1090 = \qty{3.1e10}{W}$. (b) Um terço do fluxo havaiano. (c) $B$ cresce com $A\,v$: a pluma sob a placa mais rápida deve refazer a mesma elevação duas vezes mais depressa, e seu fluxo é duas vezes maior.
\end{solution}

\begin{solution}{exo:b1:hotspots:8}
(a) $350/(16.1 - 10.3) = 350/5.8 = \qty{60}{km.Myr^{-1}}$, isto é, \qty{60}{mm.yr^{-1}}. (b) A trilha jovem, \qty{29 \pm 5}{mm.yr^{-1}}, concorda com os \qty{26 \pm 5}{mm.yr^{-1}} da placa dentro das incertezas; a parte antiga foi mais de duas vezes mais rápida, $60 - 26 = 34$, quase sete incertezas-padrão da velocidade da placa. (c) Sim: N\,\ang{54}\,L é o \omterm{def:b1:spherical-astronomy:horizontal}{azimute} \ang{54}, oposto a \ang{235} dentro de um grau. O vulcanismo se move sobre a placa no sentido oposto ao da placa sobre o \omterm{def:b1:hotspots:hotspot}{ponto quente}. (d) A distensão afastou as caldeiras mais antigas umas das outras, aumentando as distâncias entre elas e, portanto, a velocidade aparente; a própria placa não foi mais rápida.
\end{solution}

\begin{solution}{exo:b1:hotspots:9}
(a) $\num{9.2e10}/\num{e12} = \qty{0.092}{W.m^{-2}}$, \qty{92}{mW.m^{-2}}. (b) Nove a dezoito vezes o excesso observado, que pode até ser nulo. (c) $L^2/\kappa =
\num{e10}/\num{e-6} = \qty{e16}{s}$, cerca de \qty{320}{Myr}; a placa atravessa \qty{1000}{km} em $1000/92 = \qty{11}{Myr}$. O calor vindo de baixo teve tempo de atravessar apenas uma pequena parte da placa enquanto ela está sobre a elevação: a maior parte do calor da pluma permanece no \omterm{def:b1:layered-earth:layers}{manto} sob a placa e é levada embora com ela, e a superfície sente pouco dele.
\end{solution}

\begin{solution}{exo:b1:hotspots:10}
(a) Dez vezes a fusão, \qty{e7}{km^3}, de raio
$(3V/4\pi)^{1/3} = \qty{134}{km}$. (b) $\num{e7}/100 = \qty{e5}{km^2}$, um disco de $2\sqrt{\num{e5}/\pi} = \qty{357}{km}$ de diâmetro: dezoito vezes menor em área que as Armadilhas. A cabeça era maior ou fundiu menos, ou se espalhou mais larga e mais fina, ou as lavas escoaram muito além dela. (c) \qty{1}{km^3} por ano, dez vezes o suprimento atual do Kilauea. (d) Uma pluma nova deve reunir rocha quente e leve suficiente para abrir caminho para cima pelo \omterm{def:b1:layered-earth:layers}{manto} viscoso, e um grande volume sobe muito mais depressa que um pequeno; a cauda que a segue sobe pelo caminho quente e fraco que a cabeça abriu, e um conduto estreito basta.
\end{solution}

\begin{solution}{exo:b1:hotspots:11}
(a) $\qty{40}{km.Myr^{-1}}\times(81 - 47)\,\unit{Myr} = \qty{1360}{km}$. (b) A placa se move com $(-80\sin\ang{20}, 80\cos\ang{20}) = (-27.4, 75.2)$ e o \omterm{def:b1:hotspots:hotspot}{ponto quente} com $(0, -40)$: em relação ao \omterm{def:b1:hotspots:hotspot}{ponto quente}, a placa se move com $(-27.4, 115.2)$ \unit{mm.yr^{-1}}, uma velocidade de \qty{118.4}{mm.yr^{-1}} em direção a $\ang{360} - \arctan(27.4/115.2) = \ang{346.6}$. (c) A trilha daria uma velocidade \qty{48}{\percent} alta demais e um rumo \ang{6.6} para leste demais: um \omterm{def:b1:hotspots:hotspot}{ponto quente} que deriva registra o movimento da placa em relação a si mesmo, e não em relação ao \omterm{def:b1:layered-earth:layers}{manto} profundo.
\end{solution}

\begin{solution}{exo:b1:hotspots:12}
(a) $\qty{3e6}{J.m^{-3}.K^{-1}}\times\qty{1300}{K}\times\qty{e15}{m^3} =
\qty{3.9e24}{J}$. (b) Em $\qty{3.16e13}{s}$, \qty{1.2e11}{W}: cerca do calor carregado pela pluma havaiana hoje (1.3 vezes), e \qty{0.26}{\percent} dos \qty{47}{TW} da Terra. (c) Um limite inferior: a cabeça também traz o calor em excesso dos nove décimos de sua rocha que não fundem e aquece a base da \omterm{def:b1:layered-earth:lithosphere}{litosfera}, e o calor latente foi deixado de fora.
\end{solution}

\begin{solution}{pb:b1:hotspots:1}
\textbf{1.} No Havaí ($\varphi = \ang{19.0}$, $\lambda = \ang{-155.2}$),
com $\sin\ang{304} = -0.8290$ e $\cos\ang{304} = 0.5592$:
\[
\begin{aligned}
  \vec r_H &= (-0.8583,\ -0.3966,\ 0.3256), &
  \vec e &= (0.4195,\ -0.9078,\ 0),\\
  \vec n &= (0.2955,\ 0.1366,\ 0.9455), &
  \vec t_H &= (-0.1825,\ 0.8289,\ 0.5287).
\end{aligned}
\]

\textbf{2.} Em Louisville ($\varphi = \ang{-53.6}$, $\lambda = \ang{-140.6}$),
com $\sin\ang{316} = -0.6947$ e $\cos\ang{316} = 0.7193$:
\[
\begin{aligned}
  \vec r_L &= (-0.4586,\ -0.3767,\ -0.8049), &
  \vec e &= (0.6347,\ -0.7727,\ 0),\\
  \vec n &= (-0.6220,\ -0.5109,\ 0.5934), &
  \vec t_L &= (-0.8883,\ 0.1693,\ 0.4269).
\end{aligned}
\]

\textbf{3.} $\lVert\vec t_H\rVert^2 = 0.0333 + 0.6871 + 0.2795 = 0.9999$ e
$\vec t_H\cdot\vec r_H = 0.1566 - 0.3287 + 0.1721 = 0.0000$, dentro do arredondamento.

\textbf{4.} $\vec t_H\times\vec t_L = (0.2643, -0.3917, 0.7054)$, de comprimento
$0.8491$. O ângulo entre as trilhas é $\arcsin0.8491 = \ang{58.1}$ (seu produto escalar, $0.528$, é positivo, logo o ângulo é agudo).

\textbf{5.} $\vec p' = (0.3113, -0.4613, 0.8308)$. Então, pelo sinal do produto misto,
$(\vec p'\times\vec r_H)\cdot\vec t_H = -0.98 < 0$: o polo é
$\vec p = (-0.3113, 0.4613, -0.8308)$.

\textbf{6.} $\arcsin(-0.8308) = \ang{-56.18}$ e
$\operatorname{atan2}(0.4613, -0.3113) = \ang{124.01}$: o polo fica em
\ang{56.2}\,S, \ang{124.0}\,L, no Oceano Antártico, ao sul da Austrália.

\textbf{7.} $\vec p\cdot\vec r_H = 0.2672 - 0.1830 - 0.2705 = -0.1863$, logo
$\Delta_H = \ang{100.7}$; $\vec p\cdot\vec r_L = 0.1428 - 0.1738 + 0.6687 =
0.6377$, logo $\Delta_L = \ang{50.4}$.

\textbf{8.} O vetor $\vec p\times\vec r_L$ é perpendicular a $\vec p$ e a
$\vec r_L$, como $\vec t_L$ é (pela construção de $\vec p$, e porque uma trilha é horizontal); os dois estão sobre a mesma reta, e seu comprimento é
$\sin\Delta_L$. Logo $(\vec p\times\vec r_L)\cdot\vec t_L = \pm\sin\Delta_L$, e o cálculo dá $+0.770 = \sin\ang{50.4}$: a placa se move em Louisville ao longo de $+\vec t_L$, no sentido em que a trilha corre. A regra de sinal, aplicada no Havaí, poderia ter dado o sentido errado em Louisville; não deu, de modo que as duas direções são reproduzidas exatamente, \ang{304} e \ang{316}.

\textbf{9.} $\omega = 67/(6371\sin\ang{50.38}) = 67/4908 =
\qty{0.01365}{rad.Myr^{-1}}$, isto é, \qty{0.782}{\degree\per\Myr}.

\textbf{10.} $360/0.782 = \qty{460}{Myr}$.

\textbf{11.} A regra dos cossenos entre \ang{56.18}\,S, \ang{124.01}\,L e
\ang{63.78}\,S, \ang{112.48}\,L dá \ang{9.5}. A taxa das duas trilhas, \qty{0.78}{\degree\per\Myr}, fica \qty{15}{\percent} abaixo dos
\qty{0.92}{\degree\per\Myr} ajustados a treze velocidades.

\textbf{12.} Sociedade \ang{73.6}, Pitcairn \ang{77.8}, Cobb \ang{134.7}.

\textbf{13.} $\omega R = 0.01365\times6371 = \qty{86.98}{mm.yr^{-1}}$, multiplicado por
$\sin\Delta$: Sociedade \qty{83.5}{mm.yr^{-1}}, Pitcairn
\qty{85.0}{mm.yr^{-1}}, Cobb \qty{61.8}{mm.yr^{-1}}.

\textbf{14.} Pitcairn concorda: \qty{5}{mm.yr^{-1}} abaixo do medido, dentro de seus \qty{15}{mm.yr^{-1}}. Sociedade fica \qty{25.5}{mm.yr^{-1}} lenta demais, 2.6 incertezas, e Cobb fica \qty{18.8}{mm.yr^{-1}} rápida demais, mais de seis: nenhuma concorda.

\textbf{15.} O polo é perpendicular aos dois vetores das trilhas. Se eles são quase paralelos, seu produto vetorial é curto, e um erro de alguns graus em qualquer sentido o faz girar por um ângulo grande: o polo é mal determinado.

\textbf{16.} $v = \qty{92e-3}{m}/\qty{3.156e7}{s} = \qty{2.92e-9}{m.s^{-1}}$
e $B = 2300\times\num{4.83e8}\times\num{2.92e-9} = \qty{3240}{kg.s^{-1}}$.

\textbf{17.} $Q = (910/\num{3.2e-5})\times3240 = \qty{9.2e10}{W}$,
\qty{0.20}{\percent} de \qty{47}{TW}.

\textbf{18.} As dorsais produzem $21/67000 = \qty{3.1e-4}{km^3}$ por ano por
quilômetro; $0.1/\num{3.1e-4} = \qty{320}{km}$ de dorsal.

\textbf{19.} $74000/0.1 = \num{740000}$ anos.

\textbf{20.} $\num{2e6}/\num{81e6} = \qty{0.025}{km^3}$ por ano: só o Kilauea é hoje alimentado com quatro vezes a produção média de todo o \omterm{def:b1:hotspots:hotspot}{ponto quente}, que está perto de sua fase mais forte.

\textbf{21.} $v_H = 86.98\sin\ang{100.73} = 86.98\times0.9825 =
\qty{85.5}{mm.yr^{-1}}$.

\textbf{22.} $92 - 85.5 = \qty{6.5}{mm.yr^{-1}}$, 2.2 incertezas-padrão da velocidade medida: os dois \omterm{def:b1:hotspots:hotspot}{pontos quentes} não concordam bem quanto ao movimento de uma única \omterm{def:b1:plate-kinematics:rigid-plate}{placa rígida} sobre um referencial fixo.

\textbf{23.} As idades dão \qty{9.1}{cm} e \qty{8.6}{cm} por ano, médias sobre a cadeia inteira; os \qty{9.2}{cm} por ano de Morgan e Morgan ficam perto do primeiro, e a previsão de Louisville, \qty{8.5}{cm} por ano, perto do segundo. Os quatro valores ficam a menos de \qty{0.7}{cm} por ano uns dos outros: trilhas, idades e direções contam a mesma história a cerca de \qty{8}{\percent}.

\textbf{24.} Com uma casa decimal, a trilha de Louisville dá à placa do Pacífico sobre o \omterm{def:b1:hotspots:hotspot}{ponto quente} havaiano a velocidade
\textbf{\qty[text-series-to-math=true]{8.5}{cm.yr^{-1}}}.
\end{solution}
